%%%%%
%
\section{Passage d'un vecteur d'un rep�re mobile au rep\`ere g�ographique}

\itemstyle{Op\'erateur~:} \operfmt{Mobile2geo\_3D}
\index{Mobile2geo\_3D}

%%%%%OBJET%%%%
\itemstyle{Objet~:}
Cet op\'erateur calcule les composantes dans le rep�re g�ographique
local d'un vecteur 3D donn\'e dans un rep\`ere mobile 
caract\'eris\'e par 3 angles d'orientation.

%%%%%ARGUMENTS%%%%
\bigskip
%%%%%ENTREES%%%%
\begin{tabarg}
\multicolumn{3}{l}{\titretab {Entr\'ees}}\\
\hline
$X_l$ & Vecteur, \texttt{float} & Composante du vecteur d'entr\'ee suivant l'axe longitudinal de
la plateforme, positif vers l'avant \\
$X_t$ & Vecteur, \texttt{float} & Composante du vecteur d'entr\'ee suivant l'axe transversal
de la plateforme, positif vers tribord \\
$X_z$ & Vecteur, \texttt{float} & Composante du vecteur d'entr\'ee suivant 
l'axe �vertical� de la plateforme, positif vers le bas \\
$\psi$ & Vecteur, \texttt{float} & Cap de la plate-forme mobile, dans le sens 
r�trograde � partir du nord g�ographique ($degres$)\\
$\phi$ & Vecteur, \texttt{float} & Roulis de la plate-forme mobile,
positif lorsque le cot� tribord est enfonc� ($degres$)\\
$\theta$ & Vecteur, \texttt{float} & Tangage de la plate-forme mobile,
positif lorsque l'avant est relev� ($degres$)\\
\end{tabarg}

\medskip
%%%%%SORTIES%%%%
\begin{tabarg}
\multicolumn{3}{l}{\titretab {Sorties}}\\
\hline
$X_e$ & Vecteur, \texttt{float} & Composante zonale du vecteur r\'esultat \\
$X_n$ & Vecteur, \texttt{float} & Composante m\'eridienne du vecteur r\'esultat \\
$X_v$ & Vecteur, \texttt{float} & Composante verticale du vecteur r\'esultat \\
\end{tabarg}


%%%%%FORMULE%%%%
\itemstyle{Formule}


$X_e = X_l.\sin\psi.\cos\theta +$
        $X_t.(\cos\psi.\cos\phi+\sin\psi.\sin\theta.\sin\phi) +$
        $X_z.(\sin\psi.\sin\theta.\cos\phi - \cos\psi.\sin\phi)$ \\   

$X_n = X_l.\cos\psi.\cos\theta -$
        $X_t.(\sin\psi.\cos\phi - \cos\psi.\sin\theta.\sin\phi) +$
        $X_z.(\cos\psi.\sin\theta.\cos\phi + \sin\psi.\sin\phi)$\\

$X_v = X_l.\sin\theta -$ 
        $X_t.\cos\theta.\sin\phi -$
        $X_z.\cos\theta.\cos\phi$ \\



%%%%%COMMENTAIRE%%%%
\itemstyle{Commentaire~:}
Les trois angles d'attitude sont ceux dit "de Cardan" (lacet,
tangage, roulis), qui poss�dent des ind�terminations ou des non d�finitions
d'angles dans certaines configurations. Ces formules sont donc adapt�es
aux �volutions usuelles d'un bateau ou un avion, mais pas d'un mobile
totalement libre.

%%%%%AUTEUR%%%%
\itemstyle{Auteur~:}
CNRM/GMEI/TRAMM

%%%%%REFERENCE%%%%
\itemstyle{R\'ef\'erences~:}
Les angles couramment nomm�s cap, tangage, roulis  dans le domaine
a�ronautique sont pr�cis�ment d�finis dans la norme des \textit{
symboles et vocabulaire de la m�canique du vol} comme \emph{Azimut, 
assiette longitudinale, et angle de g�te}, avec les 
m�mes symboles $\Psi, \Theta, \Phi$, cf. norme AFNOR NFX02-115, page 7
\cite{AFNOR-74}.
